\documentclass[12pt]{article}

\usepackage[margin=1in]{geometry}
\usepackage{amsmath}
\usepackage{hyperref}
\usepackage{enumitem}
\usepackage{titlesec}
\usepackage{booktabs}
\usepackage{longtable}

\hypersetup{
    colorlinks=true,
    linkcolor=blue,
    citecolor=blue,
    urlcolor=blue
}

\titleformat{\section}{\large\bfseries}{\thesection}{1em}{}
\titleformat{\subsection}{\normalsize\bfseries}{\thesubsection}{1em}{}
\titleformat{\subsubsection}{\normalsize\itshape}{\thesubsubsection}{1em}{}

\title{Bills of Credit in Colonial and Revolutionary America:\\
Paper Money, Legal Tender, and Fiscal Reckoning}
\author{}
\date{}

\begin{document}
\maketitle

\tableofcontents
\newpage

\section{Introduction}

From the 1690s through the 1780s, British North American colonies and, later,
the independent American states conducted one of the most extensive experiments
with government-issued paper money---``bills of credit''---in the history of the
world before the French Revolution.  These instruments  ranged from
well-managed media of exchange that circulated near par for decades to
catastrophically depreciated scraps of paper that gave the English language the
phrase ``not worth a Continental.''  Their story illuminates fundamental
questions about money, public finance, and political economy that the Founders
confronted when they drafted the Constitution.

\section{Pre-War Colonial Precedents}

\subsection{Massachusetts and the Invention of Colonial Bills of Credit (1690)}

The first bills of credit in the British colonies were issued by Massachusetts
in December 1690 to pay soldiers returning from Sir William Phips's failed
expedition against Quebec.  The colonial treasury was empty, the troops were
mutinous, and the General Court authorized bills promising redemption from
future tax revenues.  The bills were receivable for taxes, which gave them a
measure of value.  This expedient---printing paper promises against future tax
receipts---became the template for every subsequent colonial issue.

Within a few years, Massachusetts made additional issues, and the bills began to
depreciate.  By the 1710s and 1720s, Massachusetts had a persistent pattern of
over-issue, partial depreciation, and promises of future redemption that were
repeatedly postponed.

\subsection{The Land Bank Model}

Several colonies lent bills of credit into circulation through ``loan offices''
or ``land banks.''  The mechanism was straightforward: the colonial government
would print bills and lend them to landowners on the security of their real
property, typically at 5\% interest.  As borrowers repaid principal and
interest, bills returned to the treasury and were retired---in theory.

\textbf{Pennsylvania (1723)} established the most celebrated land bank.  Its
bills circulated near par for decades, interest earnings from the loans
substantially funded the colonial government, and the system won the
enthusiastic endorsement of Benjamin Franklin, who published his \emph{A Modest
Enquiry into the Nature and Necessity of a Paper-Currency} in 1729 to advocate
expanded issues.  Even Adam Smith, no friend of paper money, conceded that
Pennsylvania's currency ``is said never to have sunk below the value of the
gold and silver which was current in the colony before the first issue of
paper money'' (\emph{Wealth of Nations}, Book~II, Chapter~2).

\textbf{New York and New Jersey} operated similar land banks with generally
moderate results, though with occasional episodes of depreciation.

\subsection{New England Excess: Rhode Island and the Inflationary Spiral}

Rhode Island became the most notorious inflator among the colonies.  Between
1710 and 1750, it made repeated large emissions that drove the bills to a
fraction of their face value.  By the 1740s, Rhode Island bills circulated in
neighboring colonies at steep discounts, creating confusion and commercial
friction throughout New England.  Connecticut and New Hampshire issued more
moderately but were dragged down by the regional currency disorders.

Massachusetts's own bills depreciated substantially: by the 1740s, the exchange
rate had deteriorated to roughly 10 or 11 old tenor shillings per 1 shilling
sterling, compared with 1.33 shillings per sterling at the official rate.

\subsection{Massachusetts Redemption of 1749}

A turning point came in 1749.  Parliament reimbursed Massachusetts
\textsterling183,649 in specie (Spanish silver dollars and copper) for its
expenses in the 1745 Louisbourg expedition.  Governor William Shirley and
Thomas Hutchinson pushed through a plan to use this windfall to redeem the
outstanding bills of credit at a rate of 10~old tenor shillings to 1 shilling
sterling and to return the colony to a hard-money standard.  This was
accomplished by 1750, making Massachusetts the first colony to voluntarily
retire its paper money and resume specie circulation.  The transition was
painful for debtors, who found their obligations suddenly revalued upward,
but it stabilized the currency.

\subsection{The Southern Colonies}

\textbf{South Carolina} issued bills of credit beginning in 1703 to finance
defense against Spanish and Indian threats.  The colony issued heavily during
Queen Anne's War and the Yamasee War (1715--17), and its bills depreciated to
roughly 7 South Carolina pounds per 1 pound sterling by the 1720s.  After
mid-century, the currency stabilized at approximately that ratio and circulated
with reasonable predictability.

\textbf{Virginia} was a relative latecomer to paper money, issuing its first
bills of credit in 1755 to finance its role in the French and Indian War.  It
authorized approximately \textsterling250,000 in bills between 1755 and 1762.
These were made legal tender and promised redemption from tax revenues.  They
depreciated modestly but remained in tolerable circulation.

\textbf{North Carolina, Georgia, and Maryland} each issued bills with varying
degrees of depreciation and controversy.

\section{British Restrictions: The Currency Acts of 1751 and 1764}

\subsection{The Currency Act of 1751}

In response to the chronic depreciation of New England currencies and the
complaints of British merchants who were paid in depreciated colonial paper,
Parliament passed the \textbf{Currency Act of 1751}.  It applied only to the
four New England colonies (Massachusetts, Connecticut, Rhode Island, and
New Hampshire) and provided that:

\begin{itemize}[nosep]
    \item No new bills of credit could be made legal tender for private debts.
    \item Bills issued for current government expenses must specify a
          redemption date no more than two years in the future and must be
          funded by taxes levied within that period.
    \item Existing legal tender provisions would expire as the bills were
          redeemed.
\end{itemize}

Since Massachusetts had already retired its bills in 1749--50, the Act
principally constrained Rhode Island, which bitterly resisted and continued to
find ways to evade its spirit.

\subsection{The Currency Act of 1764}

The more consequential statute was the \textbf{Currency Act of 1764}, which
extended the prohibition on legal tender bills of credit to \emph{all} the
colonies.  Its immediate trigger was Virginia's wartime issues: British
merchants trading in Virginia complained that they were compelled to accept
depreciated paper at face value under Virginia's legal tender law.

The Act provided that:

\begin{itemize}[nosep]
    \item No colony could make new or existing bills of credit legal tender in
          payment of any debts.
    \item Bills already in circulation could continue until their specified
          retirement dates but could not be re-issued or extended.
    \item Governors who assented to acts violating these provisions would be
          dismissed.
\end{itemize}

The 1764 Act became a significant colonial grievance.  It constrained the money
supply in colonies that had managed their currencies responsibly
(Pennsylvania in particular protested vigorously) as well as those that had
not.  Benjamin Franklin, appearing before Parliament in 1766, argued that the
restriction on paper money was as important a source of colonial discontent as
the Stamp Act itself.

\subsection{Partial Relaxation (1770 and 1773)}

Under persistent colonial lobbying, Parliament partially relented.  The
\textbf{1770 amendment} allowed New York to issue bills receivable for public
dues (taxes and fees) but not legal tender for private debts.  The \textbf{1773
amendment} extended this permission to all colonies: they could issue bills
that were legal tender for \emph{public} payments (taxes, fees, fines) but not
for private debts.  This compromise satisfied some colonies but not others, and
the underlying grievance persisted until independence made the question moot.

\section{Adam Smith on Colonial Paper Money}

Adam Smith devotes considerable attention to colonial bills of credit in
\emph{The Wealth of Nations} (1776), principally in Book~II, Chapter~2 and in
Book~IV (on colonies).  His analysis is characteristically nuanced:

\begin{enumerate}
    \item \textbf{Pennsylvania as a success:} Smith acknowledges that
    Pennsylvania's bills ``never sunk below the value of the gold and silver
    which was current in the colony before the first issue of paper money.''
    He attributes this to moderate quantities, prudent management, and the
    colony's prosperous economy.

    \item \textbf{New England as a failure:} Smith treats the New England
    currencies, especially those of Rhode Island and Massachusetts before 1749,
    as cautionary examples.  He notes that the bills depreciated because the
    colonies issued more than the economy could absorb and failed to redeem
    them on schedule.

    \item \textbf{Legal tender as the crucial vice:} Smith argues that making
    paper money legal tender for private debts is the decisive error.  It
    allows the government to force creditors to accept depreciated paper at
    face value, which is ``a sort of injustice'' and ``in reality a scheme of
    fraudulent debtors to cheat their creditors'' (Book~II, Chapter~2).

    \item \textbf{The 1764 Act:} Smith expresses qualified support for the
    prohibition on legal tender but criticizes the blanket nature of the
    restriction, suggesting that colonies which had managed their paper
    responsibly might have been permitted to continue issuing non-legal-tender
    bills receivable for taxes.
\end{enumerate}

\section{Paper Money During the War of Independence}

\subsection{The Continental Currency}

Once war began, both the Continental Congress and the individual states resorted
to paper money on a vast scale.  The Continental Congress, lacking the power to
tax, had few alternatives.

\begin{itemize}
    \item \textbf{Total authorized issues:} Between June 1775 and November 1779,
    Congress authorized approximately \$241.6~million (face value) in
    Continental currency.  The bills were denominated in ``Spanish milled
    dollars'' and initially bore promises that the ``United Colonies'' would
    redeem them, with the understanding that the states would levy taxes to
    fund their retirement.

    \item \textbf{No taxing power:} Congress could not tax.  It could only
    ``requisition'' the states to collect taxes and remit the proceeds.  The
    states chronically under-complied.  Without a reliable revenue stream
    backing them, the Continentals depended entirely on confidence---which
    eroded with each new issue.

    \item \textbf{Depreciation:} The Continentals held near par through late
    1776.  By early 1778, they had fallen to roughly 4-to-1 against specie.
    By early 1779, the ratio was approximately 10-to-1.  By the end of 1779,
    it was 40-to-1.  By early 1781, Continental bills were trading at
    roughly 100-to-1 or more, and they effectively ceased to circulate.

    \item \textbf{The ``Forty-for-One'' scheme (March 1780):} In a pivotal act,
    Congress acknowledged the depreciation and resolved that old Continentals
    could be exchanged at 40-to-1 for new bills backed by state revenues.  This
    amounted to an official devaluation of 97.5\%.  Even the new bills
    subsequently depreciated.
\end{itemize}

\subsection{State Bills of Credit During the War}

In addition to the Continentals, individual states issued their own bills of
credit to finance war expenditures.  Rothbard and other authorities document
these extensively.

\begin{itemize}
    \item \textbf{Total state issues:} The thirteen states collectively issued
    approximately \$209~million (face value) in state bills of credit between
    1775 and 1781.

    \item \textbf{Virginia} was the largest single issuer among the states,
    authorizing over \$128~million in bills between 1775 and 1781.
    Virginia's bills depreciated roughly in tandem with the Continentals.

    \item \textbf{The Carolinas:} South Carolina and North Carolina each issued
    substantial quantities.  South Carolina's paper depreciated significantly
    but was eventually partially redeemed.  North Carolina's bills fell to a
    small fraction of face value.

    \item \textbf{New England states:} Massachusetts, Connecticut, and others
    issued more moderately but still experienced depreciation, though generally
    less extreme than the Southern states.

    \item \textbf{Pennsylvania, New York, New Jersey, Maryland:} Issued moderate
    amounts.  Pennsylvania's bills, building on its colonial tradition,
    held their value relatively better than those of most other states, though
    they too depreciated as the war dragged on.

    \item \textbf{Georgia:} Issued heavily relative to its small economy; its
    bills depreciated drastically.
\end{itemize}

\subsection{Legal Tender Provisions}

Both the Continental currency and the state bills of credit were typically
declared \textbf{legal tender}.  This was the very practice that the British
Currency Acts of 1751 and 1764 had prohibited and that Adam Smith had denounced.

\begin{itemize}
    \item \textbf{Continental Congress resolutions:} In January 1776, Congress
    resolved that anyone who refused to accept Continental currency at face
    value should ``be deemed, published, and treated as an enemy of his
    country, and precluded from all trade or intercourse with the inhabitants
    of these colonies.''  This was social and economic coercion, though Congress
    lacked formal enforcement machinery.

    \item \textbf{State legal tender laws:} Most states enacted statutes making
    both Continentals and their own state bills legal tender for all debts,
    public and private.  Penalties for refusing them varied:
    \begin{itemize}[nosep]
        \item Fines (often equal to the value of the debt).
        \item Forfeiture of the debt owed (i.e., the debtor was discharged).
        \item Being declared ``an enemy of the state.''
        \item In some cases, imprisonment.
    \end{itemize}

    \item \textbf{Price controls:} Several states attempted to support the
    currency through price-fixing legislation (the ``regulation'' controversy
    of 1776--1778).  New England states held conventions in 1776--77 to
    coordinate price ceilings.  The Mid-Atlantic states attempted similar
    schemes.  These efforts uniformly failed; goods disappeared from legal
    markets, and the price controls were abandoned by 1778--79.

    \item \textbf{Tender laws and injustice:} The legal tender provisions
    generated enormous resentment among creditors, who were forced to accept
    rapidly depreciating paper at face value.  As Pelatiah Webster, a
    Philadelphia merchant and pamphleteer, argued in 1780, the tender laws
    effected a massive and arbitrary redistribution of wealth from creditors to
    debtors---precisely the ``fraud'' that Adam Smith had identified.
\end{itemize}

\subsection{What Happened to These Debts?}

The fate of the various paper currencies after the war divides into several
categories:

\subsubsection{Continental Currency}

The Continentals were never redeemed at anything close to face value.  The
sequence was:

\begin{enumerate}[nosep]
    \item The 40-to-1 conversion of March 1780 implicitly repudiated 97.5\%
          of the face value.
    \item Even the ``new tenor'' bills issued under the 40-to-1 scheme
          subsequently depreciated.
    \item Under Hamilton's funding program (the \textbf{Funding Act of 1790}),
          holders of old Continental currency could exchange it for new federal
          bonds at a rate of \textbf{100-to-1}---that is, \$100 in old
          Continentals yielded \$1 in new 6\% federal bonds.  This was a
          redemption of sorts, but at one cent on the dollar.
    \item Much of the Continental currency had already been lost, destroyed, or
          abandoned.  A significant fraction was never presented for
          conversion.  The total amount actually exchanged under the 1790 Act
          was far less than the nominal \$200+ million originally issued.
\end{enumerate}

\subsubsection{State Bills of Credit}

The states dealt with their wartime paper in diverse ways:

\begin{itemize}
    \item \textbf{Redemption at depreciated rates:} Most states that redeemed
    their bills did so at rates reflecting the actual depreciation.  Virginia,
    for example, scaled its wartime bills at various ratios (ultimately as
    steep as 1000-to-1 for the most depreciated issues) and converted the
    remainder into state bonds or certificates.

    \item \textbf{Tax absorption:} Some states declared their bills receivable
    for taxes at stated ratios and gradually retired them as taxes were
    collected.  This was the orderly method, but it presupposed that the state
    could actually collect taxes---a significant challenge in the
    war-disrupted economy.

    \item \textbf{Effective repudiation:} Several states allowed their bills to
    become effectively worthless without any formal act of repudiation.
    Georgia's bills, for instance, depreciated to the vanishing point.
    North Carolina's wartime paper similarly lost virtually all value.

    \item \textbf{Federal assumption:} Under Hamilton's assumption plan
    (\textbf{Assumption Act of 1790}), the federal government assumed
    approximately \$25~million in state debts.  However, the debts assumed
    were predominantly state \emph{bonds and certificates} (funded debts
    owed to soldiers, suppliers, and creditors), not the bills of credit
    themselves, most of which had already been retired, repudiated, or
    converted before 1790.
\end{itemize}

\subsubsection{The Post-War Paper Money Crisis (1785--1787)}

After the war, several states resumed issuing bills of credit, reigniting the
controversy:

\begin{itemize}
    \item \textbf{Rhode Island (1786)} passed a particularly aggressive paper
    money law, issuing \textsterling100,000 in legal tender bills and imposing
    draconian penalties on anyone who refused them.  The resulting legal
    conflict produced the famous case of \emph{Trevett v.\ Weeden} (1786),
    in which a Rhode Island court effectively nullified the tender law.

    \item \textbf{Seven states} issued new paper money between 1785 and 1787:
    Rhode Island, New York, New Jersey, Pennsylvania, the Carolinas, and
    Georgia.

    \item These post-war issues were a direct impetus for the Constitutional
    Convention.  The perceived abuses of state paper money motivated the
    provisions in \textbf{Article~I, Section~10} of the Constitution:
    ``No State shall \ldots\ emit Bills of Credit; \ldots\ make any Thing but
    gold and silver Coin a Tender in Payment of Debts.''
\end{itemize}

\section{Summary: Patterns and Legacies}

Several patterns emerge from this extended experiment with paper money:

\begin{enumerate}
    \item \textbf{Moderate issues could succeed.} Pennsylvania's colonial land
    bank demonstrated that paper money, issued in limited quantities against
    real assets and not made legal tender for private debts, could circulate
    near par and even benefit the public finances.  The key ingredients were
    restraint in quantity, reliable mechanisms for retirement, and a productive
    economy.

    \item \textbf{Legal tender was the fatal temptation.} When governments
    made depreciating paper legal tender, they transferred wealth from
    creditors to debtors by fiat, undermined commercial confidence, and
    created incentives for hoarding specie and evading the law.  Every colony
    and state that made its bills legal tender for private debts experienced
    these pathologies.

    \item \textbf{Wartime finance overwhelmed restraint.}  The Revolutionary
    War created fiscal pressures that no colonial precedent had contemplated.
    The Continental Congress and the states printed money because they could
    not tax effectively, and the resulting inflation was a de facto tax---one
    that fell erratically and unjustly.

    \item \textbf{Depreciation was a form of taxation and repudiation.}
    As the Continentals fell from par to one-hundredth of their face value,
    each holder bore a loss proportional to his holdings.  The process was,
    as contemporaries recognized, a tax on holders of money---but an
    arbitrary one, unlegislated and unpredictable.

    \item \textbf{The British restrictions had a point.}  The Currency Acts of
    1751 and 1764, resented as imperial overreach, addressed a real problem:
    colonial legislatures had repeatedly succumbed to the temptation to
    over-issue and to shift losses to creditors through legal tender laws.
    The revolutionary experience vindicated the British concern even as it
    rendered the British remedy irrelevant.

    \item \textbf{The Constitutional prohibition was born of experience.}
    The Framers, having witnessed the Continentals' collapse and the post-war
    state paper money abuses, wrote the prohibition on state bills of credit
    and state legal tender laws directly into the Constitution.  This was not
    abstract theorizing; it was a considered response to a half-century of
    evidence.
\end{enumerate}

\section{A Note on Sources}

Murray Rothbard's \emph{Conceived in Liberty}, Volume~4 (\emph{The
Revolutionary War, 1775--1784}), provides a detailed libertarian critique of
the revolutionary paper money experiment, emphasizing the coercive character
of legal tender laws and price controls.  Other important sources include:

\begin{itemize}[nosep]
    \item Adam Smith, \emph{The Wealth of Nations} (1776), Book~II, Chapter~2
          and Book~IV, Chapter~7, Part~2.
    \item E.\ James Ferguson, \emph{The Power of the Purse: A History of
          American Public Finance, 1776--1790} (1961).
    \item Leslie V.\ Brock, \emph{The Currency of the American Colonies,
          1700--1764} (1975).
    \item Ben Baack, ``Forging a Nation State: The Continental Congress and the
          Financing of the War of American Independence,'' \emph{Economic
          History Review} (2001).
    \item Farley Grubb, ``State Redemption of the Continental Dollar,
          1779--90,'' \emph{William and Mary Quarterly} (2012), and numerous
          other articles on colonial and revolutionary currency.
    \item Charles W.\ Calomiris, ``Institutional Failure, Monetary Scarcity,
          and the Depreciation of the Continental,'' \emph{Journal of Economic
          History} (1988).
    \item Ron Michener, ``Backing Theories and the Currencies of
          Eighteenth-Century America,'' \emph{Journal of Economic History}
          (1988).
    \item Pelatiah Webster, \emph{Political Essays on the Nature and Operation
          of Money, Public Finances, and Other Subjects} (1791).
\end{itemize}

\end{document}
